% generator_I48.tex -- Prop I.48 (Theor.): converse of Pythagorean theorem.
\documentclass[12pt]{article}\usepackage{geometry}\geometry{a5paper,margin=1.5cm}
\usepackage[lining]{ebgaramond}\usepackage{amssymb}\usepackage{byrne}
\def\mpPre{u := 1cm; textLabels := false;}
\begin{document}
\defineNewPicture[1/4]{%
  pair A,B,C,D;numeric d;d:=7/4u;%
  A:=(0,0);B:=A shifted (0,4u);C:=A shifted (d,0);D:=A shifted (-d,0);%
  byAngleDefine(B,A,C,byred,SOLID_SECTOR);%
  byAngleDefine(D,A,B,byyellow,SOLID_SECTOR);%
  draw byNamedAngleResized();%
  draw byLine(A,B,byblue,SOLID_LINE,REGULAR_WIDTH);%
  byLineDefine(A,C,byblack,SOLID_LINE,REGULAR_WIDTH);%
  byLineDefine(A,D,byblack,DASHED_LINE,REGULAR_WIDTH);%
  byLineDefine(B,C,byred,SOLID_LINE,REGULAR_WIDTH);%
  byLineDefine(B,D,byred,DASHED_LINE,REGULAR_WIDTH);%
  draw byNamedLineSeq(0)(AC,AD,BD,BC);%
  draw byLabelsOnPolygon(D,B,C,A)(ALL_LABELS,0);%
}
\noindent\begin{minipage}[t]{0.45\linewidth}\centering\vspace{0pt}%
\drawCurrentPicture\end{minipage}\hfill
\begin{minipage}[t]{0.52\linewidth}\vspace{0pt}%
\textbf{Book~I.\enspace Prop.\ XLVIII. Theor.}
\medskip\noindent\itshape
If the square of one side (\drawUnitLine{BC}) of a triangle is
equal to the squares of the other two sides (\drawUnitLine{AB}
and \drawUnitLine{AC}), the angle (\drawAngle{BAC}) subtended
by that side is a right angle.
\bigskip\noindent\upshape\begin{center}
Draw $\drawUnitLine{AD}\perp\drawUnitLine{AB}$ and
$=\drawUnitLine{AC}$~(pr.~I.11,~I.3)\\
and draw \drawUnitLine{BD} also.\\[6pt]
Since $\drawUnitLine{AD}=\drawUnitLine{AC}$~(const.)\\
$\drawUnitLine{AD}^2=\drawUnitLine{AC}^2$;\\[4pt]
$\therefore \drawUnitLine{AD}^2+\drawUnitLine{AB}^2=\drawUnitLine{AC}^2+\drawUnitLine{AB}^2$\\[4pt]
but $\drawUnitLine{AD}^2+\drawUnitLine{AB}^2=\drawUnitLine{BD}^2$~(pr.~I.47),\\
and $\drawUnitLine{AC}^2+\drawUnitLine{AB}^2=\drawUnitLine{BC}^2$~(hyp.)\\[4pt]
$\therefore \drawUnitLine{BD}^2=\drawUnitLine{BC}^2$,\\[4pt]
$\therefore \drawUnitLine{BD}=\drawUnitLine{BC}$;\\[4pt]
and $\therefore \drawAngle{DAB}=\drawAngle{BAC}$~(pr.~I.8),\\[4pt]
consequently \drawAngle{BAC} is a right angle.
\end{center}
\begin{flushright}Q.\ E.\ D.\end{flushright}
\end{minipage}
\end{document}
